
\documentclass[10pt,conference,compsoc]{IEEEtran}

% --- Packages ---
\usepackage[utf8]{inputenc}
\usepackage[T1]{fontenc}
\usepackage[french]{babel}
\usepackage{amsmath, amssymb, amsthm}
\usepackage{amsfonts}
\usepackage{xcolor}
\usepackage{booktabs}
\usepackage{graphicx}
\usepackage{listings}
\usepackage{array}
\usepackage{url}
\usepackage{microtype}
\usepackage{multirow}
\usepackage{siunitx}

% --- Configuration listings ---
\lstset{
    language=C++,
    basicstyle=\ttfamily\scriptsize,
    keywordstyle=\color{blue}\bfseries,
    commentstyle=\color{gray}\itshape,
    stringstyle=\color{purple},
    numbers=left,
    numberstyle=\tiny\color{gray},
    breaklines=true,
    frame=tb,
    captionpos=b,
    tabsize=4
}

\begin{document}

\title{Télémétrie Temporelle Différentielle en Espace Utilisateur~:
Caractérisation Distante et Limites de Résilience du \textit{Localhost Security Gap}}

\author{\IEEEauthorblockN{Francois Lionnel Bala Andegue}
\IEEEauthorblockA{Département de Génie Informatique et Cybersécurité\\
École Nationale Supérieure Polytechnique de Yaoundé (ENSPY)\\
Yaoundé, Cameroun\\
Email : balaandeguefrancoislionnel@gmail.com}
}

\maketitle

% =========================================================================
\begin{abstract}
L'exposition de serveurs de développement sur un réseau local (LAN)
via une configuration d'écoute permissive (\texttt{0.0.0.0}) constitue
un vecteur d'attaque systématiquement sous-estimé dans les environnements
institutionnels. Cet article formalise la notion de \textit{Localhost
Security Gap} (LSG) distant et propose un cadre mathématique pour
l'isoler et le quantifier à distance. Nous présentons un framework
d'évaluation en espace utilisateur, développé en C++ avec la bibliothèque
\texttt{libcurl}, qui exploite la décomposition du temps de cycle HTTP
afin d'isoler le temps de traitement applicatif pur
($T_{\mathrm{calcul}}$) de la gigue réseau ($\varepsilon$). Un mécanisme
de double sonde (\textit{Dual-Probe Verification}) permet de confirmer
statistiquement toute dégradation détectée. Les expériences menées sur
un banc physique LAN (IEEE~802.3) avec quatre serveurs de développement
représentatifs (Node.js, Express, PHP intégré, Python \texttt{http.server})
démontrent que la télémétrie passive permet de discriminer le profil de
charge computationnelle de chaque serveur et de détecter des états de
saturation avec un z-score $|z| > 3$ ($p < 0{,}001$). Ces résultats
constituent une base rigoureuse pour la sensibilisation au durcissement
des configurations en environnement académique et institutionnel.
\end{abstract}

\begin{IEEEkeywords}
Télémétrie réseau, canaux auxiliaires temporels, Localhost Security Gap,
analyse de latence applicative, C++, résilience, durcissement, LAN.
\end{IEEEkeywords}

% =========================================================================
\section{Introduction}
\label{sec:introduction}
% =========================================================================

Durant le cycle de développement logiciel, les ingénieurs et étudiants
déploient massivement des serveurs web intégrés pour le débogage et les
tests. Ces services, conçus à l'origine pour une exécution isolée sur
l'interface de boucle locale (\texttt{127.0.0.1}), sont fréquemment
configurés pour écouter sur l'interface générique
\texttt{INADDR\_ANY}~(\texttt{0.0.0.0}) afin de faciliter les tests
depuis d'autres postes. Cette déviation expose l'ensemble du segment
LAN à des services non durcis~: absence de chiffrement TLS, pas de
mécanisme d'authentification, et exécution en mode mono-thread incapable
d'absorber une charge concurrente.

La littérature de sécurité adresse largement les architectures de
production (pare-feux applicatifs, load balancers), mais délaisse la
caractérisation systématique de ces environnements de développement
exposés. La question de savoir si un observateur passif sur le même
LAN peut \textit{identifier} et \textit{quantifier} la faiblesse
résiduelle de ces serveurs via la seule télémétrie HTTP reste ouverte.

Les \textbf{contributions principales} de cet article sont les suivantes~:
\begin{enumerate}
\item Un modèle mathématique formel pour l'isolation du temps de
  traitement applicatif ($T_{\mathrm{calcul}}$) à partir du Round-Trip
  Time, incluant un estimateur robuste basé sur le filtrage IQR.
\item L'implémentation d'un framework d'audit en espace utilisateur
  (C++17/\texttt{libcurl}), ne requérant aucun privilège noyau et
  intégrant une restriction de sécurité RFC~1918 pour prévenir tout
  détournement malveillant.
\item Un mécanisme de vérification par double sonde
  (\textit{Dual-Probe}) qui distingue les vraies dégradations
  applicatives des artefacts de gigue réseau.
\item Une évaluation expérimentale sur banc physique comparant quatre
  types de serveurs de développement, avec analyse statistique
  (z-score de Welch, effect size de Cohen).
\end{enumerate}

Le reste de l'article est organisé comme suit~: la
Section~\ref{sec:etat_art} présente les travaux connexes~;
la Section~\ref{sec:modelisation} développe le modèle mathématique~;
la Section~\ref{sec:architecture} décrit l'architecture du framework~;
la Section~\ref{sec:resultats} présente et analyse les résultats
expérimentaux~; la Section~\ref{sec:ethique} traite des aspects
éthiques~; la Section~\ref{sec:conclusion} conclut.

% =========================================================================
\section{État de l'Art et Travaux Connexes}
\label{sec:etat_art}
% =========================================================================

\subsection{Canaux auxiliaires temporels sur réseau}
L'exploitation des canaux auxiliaires temporels (\textit{Timing
Side-Channels}) pour déduire l'état interne d'un système distant a été
formalisée par Brumley et Boneh~\cite{brumley2003remote}, qui ont
démontré que de faibles variations de latence TCP/IP suffisent à extraire
des clés RSA d'un serveur OpenSSL distant. Ces travaux établissent le
fondement théorique de notre approche~: si des variations de l'ordre de
quelques microsecondes permettent de déduire des secrets cryptographiques,
des variations millimétriques suffisent à caractériser le profil
computationnel d'un serveur de développement.

\subsection{Empreinte de piles réseau}
L'analyse de l'empreinte des piles réseau distantes (\textit{OS
Fingerprinting}) a été formalisée par Fyodor~\cite{fyodor1998} via des
analyses de paquets TCP/IP. Notre approche opère au niveau applicatif
(HTTP), ce qui la distingue des méthodes d'empreinte de pile TCP/IP~:
elle est réalisable en espace utilisateur pur, sans forger de paquets
bruts.

\subsection{Déni de service lent sur serveurs HTTP non durcis}
Cambiaso et al.~\cite{cambiaso2013slow} ont catégorisé les attaques de
déni de service à bas débit (\textit{Slow DoS}), démontrant que les
serveurs HTTP mono-thread subissent une saturation par épuisement de
leur unique socket d'acceptation sous une charge minimale. Notre
framework mesure précisément ce phénomène via la dégradation de
$T_{\mathrm{calcul}}$ sous charge concurrente.

\subsection{Surface d'exposition des services localhost}
La problématique de l'exposition non intentionnelle de services sur
\texttt{0.0.0.0} a été analysée dans le contexte des navigateurs par
Barth et al.~\cite{barth2008robust}, qui identifient des vecteurs
d'attaque par rebond (\textit{DNS rebinding}) exploitant précisément
cette confusion entre liaison locale et liaison générique.

% =========================================================================
\section{Modélisation Mathématique de la Télémétrie}
\label{sec:modelisation}
% =========================================================================

\subsection{Décomposition du Round-Trip Time}

Soient une machine d'audit $A$ et une machine cible $B$ reliées sur un
même segment LAN. $A$ émet une requête HTTP sur le port $p \in \mathcal{P}$
de $B$, avec $\mathcal{P} = \{3000, 5000, 8080, 8082\}$. La mesure
retournée par $A$, notée $T_{\mathrm{total}}$, se décompose en~:
\begin{equation}
    T_{\mathrm{total}}(p) = T_{\mathrm{réseau}} + T_{\mathrm{calcul}}(p)
    + \varepsilon
    \label{eq:decomp}
\end{equation}
où $T_{\mathrm{réseau}} \in \mathbb{R}^+$ est le temps de transit aller-retour
des paquets sur le support physique~;
$T_{\mathrm{calcul}}(p) \in \mathbb{R}^+$ est le temps de traitement
applicatif du serveur sur le port $p$~; et $\varepsilon$ est la gigue
réseau (\textit{jitter}), modélisée comme une variable aléatoire de
distribution approximativement normale en environnement LAN contrôlé~:
$\varepsilon \sim \mathcal{N}(0, \sigma_{\varepsilon}^2)$.

\subsection{Isolation du temps de traitement applicatif}

La bibliothèque \texttt{libcurl} expose deux métriques internes qui
permettent d'affiner la décomposition~:
\begin{itemize}
    \item $t_{\mathrm{connect}}$~: durée de l'établissement de la
      connexion TCP (poignée de main à 3 voies)~;
    \item $t_{\mathrm{TTFB}}$ (Time To First Byte, \texttt{CURLINFO\_STARTTRANSFER\_TIME})~:
      durée jusqu'à la réception du premier octet de réponse.
\end{itemize}
La grandeur d'intérêt — le temps de traitement serveur pur — est alors~:
\begin{equation}
    \hat{T}_{\mathrm{calcul}} = t_{\mathrm{TTFB}} - t_{\mathrm{connect}}
    \label{eq:tcalcul}
\end{equation}
Cette estimation soustrait la latence réseau de bas niveau et isole la
phase pendant laquelle le serveur traite la requête avant d'envoyer
le premier octet de réponse.

\subsection{Calibration et estimateur robuste}

Lors d'une phase de calibration à charge minimale ($N$ requêtes légères
en séquence), $A$ constitue un vecteur de mesures
$\mathbf{x} = (x_1, \ldots, x_N)$, $x_i = \hat{T}_{\mathrm{calcul}}^{(i)}$.
Pour réduire l'influence des valeurs aberrantes (pics dus à
l'ordonnancement système, cache miss, collecte de mémoire), nous
appliquons un filtre IQR~:
\begin{equation}
    \mathcal{X}_{\mathrm{filtré}} = \left\{ x_i \;\middle|\;
    Q_1 - 1{,}5 \cdot \mathrm{IQR} \leq x_i \leq Q_3 + 1{,}5 \cdot \mathrm{IQR} \right\}
    \label{eq:iqr}
\end{equation}
L'estimateur robuste de référence est~:
\begin{equation}
    \bar{T}_{\mathrm{base}} = \frac{1}{|\mathcal{X}_{\mathrm{filtré}}|}
    \sum_{x_i \in \mathcal{X}_{\mathrm{filtré}}} x_i
    \label{eq:baseline}
\end{equation}
Toute déviation sous charge est quantifiée par~:
\begin{equation}
    \Delta T(p) = \bar{T}_{\mathrm{load}}(p) - \bar{T}_{\mathrm{base}}(p)
    \label{eq:delta}
\end{equation}

\subsection{Test statistique de dégradation}

Pour distinguer une vraie dégradation applicative du bruit $\varepsilon$,
nous formulons un test d'hypothèse~:
\begin{align}
    H_0 &: \Delta T(p) = 0 \quad \text{(serveur non dégradé)} \\
    H_1 &: \Delta T(p) > k\,\hat{\sigma} \quad \text{(dégradation significative)}
\end{align}
Le z-score de Welch~\eqref{eq:welch} est calculé pour chaque port~:
\begin{equation}
    z_p = \frac{\bar{T}_{\mathrm{load}}(p) - \bar{T}_{\mathrm{base}}(p)}
    {\hat{\sigma}_{\mathrm{base}} / \sqrt{n_{\mathrm{load}}}}
    \label{eq:welch}
\end{equation}
$H_0$ est rejetée si $|z_p| > 3$, correspondant à $p < 0{,}001$.

% =========================================================================
\section{Architecture Technique du Framework}
\label{sec:architecture}
% =========================================================================

Le framework s'exécute exclusivement en espace utilisateur. Son
architecture est segmentée en quatre modules~: (1)~validateur d'adresses
RFC~1918, (2)~moteur de sonde \texttt{libcurl}, (3)~pipeline statistique
IQR, et (4)~mécanisme de double sonde.

\subsection{Moteur de sonde et isolation temporelle}

Le Listing~\ref{lst:probe} illustre la fonction de sonde centrale. La
mesure utilise \texttt{CURLINFO\_STARTTRANSFER\_TIME} (TTFB) plutôt que
\texttt{CURLINFO\_TOTAL\_TIME}~: cette distinction est fondamentale car
elle exclut le temps de transfert du corps de la réponse (réseau pur)
et ne conserve que le délai de traitement serveur, conformément à
l'équation~\eqref{eq:tcalcul}.

\begin{lstlisting}[caption={Moteur d'acquisition temporelle — isolation du $T_{\mathrm{calcul}}$},
                   label={lst:probe}]
Measurement single_probe(const std::string& url,
                          long timeout_ms, const char* post)
{
    Measurement m;
    CURL* h = curl_easy_init();

    curl_easy_setopt(h, CURLOPT_URL, url.c_str());
    curl_easy_setopt(h, CURLOPT_WRITEFUNCTION, _discard);
    curl_easy_setopt(h, CURLOPT_TIMEOUT_MS, timeout_ms);
    curl_easy_setopt(h, CURLOPT_NOSIGNAL, 1L);
    if (post)
        curl_easy_setopt(h, CURLOPT_POSTFIELDS, post);

    CURLcode rc = curl_easy_perform(h);
    if (rc == CURLE_OK) {
        double ct, st;
        curl_easy_getinfo(h, CURLINFO_CONNECT_TIME, &ct);
        // TTFB : temps jusqu'au 1er octet de reponse
        curl_easy_getinfo(h,
            CURLINFO_STARTTRANSFER_TIME, &st);
        m.connect_ms     = ct * 1e3;
        m.ttfb_ms        = st * 1e3;
        // T_calcul = TTFB - temps de poignee de main TCP
        m.app_latency_ms = (st - ct) * 1e3;
        m.success = true;
    }
    curl_easy_cleanup(h);
    return m;
}
\end{lstlisting}

\subsection{Filtre IQR et pipeline statistique}

Le filtre IQR (équation~\eqref{eq:iqr}) est appliqué à chaque vecteur
d'échantillons avant le calcul de $\bar{T}_{\mathrm{base}}$. Le
Listing~\ref{lst:iqr} en montre l'implémentation~:

\begin{lstlisting}[caption={Filtre IQR et calcul des statistiques descriptives},
                   label={lst:iqr}]
Stats compute_stats(std::vector<double> v) {
    Stats s;
    if (v.size() < 4) return s;
    std::sort(v.begin(), v.end());

    double q1  = v[v.size() / 4];
    double q3  = v[3 * v.size() / 4];
    double iqr = q3 - q1;
    // Suppression des valeurs aberrantes (Tukey x1.5)
    v.erase(std::remove_if(v.begin(), v.end(),
        [&](double x){
            return x < q1 - 1.5*iqr || x > q3 + 1.5*iqr;
        }), v.end());

    s.count = v.size();
    s.mean  = std::accumulate(v.begin(), v.end(), 0.0)
              / s.count;
    // Calcul ecart-type, mediane, P95, P99...
    return s;
}
\end{lstlisting}

\subsection{Mécanisme de double sonde (\textit{Dual-Probe})}

Lorsqu'un seuil $\Delta T > 3\hat{\sigma}$ est détecté sur une requête
complexe (POST avec charge utile), une requête témoin (GET sur la
racine, réponse statique) est immédiatement émise. Le ratio
de corrélation~:
\begin{equation}
    r = \frac{\Delta T_{\mathrm{complexe}}}{\Delta T_{\mathrm{statique}}}
    \label{eq:ratio}
\end{equation}
permet de distinguer deux cas~:
\begin{itemize}
    \item $r \approx 1$~: les deux requêtes sont affectées de la même
      façon~; la cause est réseau (pic de gigue) → $H_0$ conservée.
    \item $r \gg 1$ (typiquement $r > 2$)~: seule la requête complexe
      est dégradée~; la cause est applicative → $H_1$ confirmée.
\end{itemize}

% =========================================================================
\section{Protocole Expérimental et Résultats}
\label{sec:resultats}
% =========================================================================

\subsection{Configuration du banc d'essai}

Le protocole repose sur deux machines physiques reliées par un commutateur
Ethernet (IEEE~802.3, \SI{100}{\mega\bit\per\second})~:
\begin{itemize}
    \item \textbf{Machine d'audit~$A$~:} Intel Core~i7, 16~Go RAM,
      Linux kernel~6.x, compilateur GCC~12.
    \item \textbf{Machine cible~$B$~:} processeur bureautique standard,
      8~Go RAM, exécutant simultanément quatre serveurs de développement
      représentatifs~: Node.js HTTP natif (port~3000), Node.js Express
      (port~5000), PHP built-in (port~8082), Python~\texttt{http.server}
      (port~8080).
\end{itemize}

Les quatre serveurs implémentent un endpoint léger (\texttt{GET /}) et
un endpoint lourd (\texttt{POST /api/compute}) qui exécute un hachage
SHA-256 itératif sur la charge utile reçue. Ce double endpoint permet
de mesurer séparément $T_{\mathrm{calcul}}$ en régime nominal et sous
charge computationnelle.

Le framework est paramétré avec $N = 100$ requêtes de calibration et
$M = 60$ requêtes de mesure par port. Le stress test utilise
$C \in \{5, 10, 20, 50\}$ requêtes simultanées pour caractériser la
courbe de saturation.

\subsection{Résultats de télémétrie différentielle}

% -----------------------------------------------------------------------
% INSTRUCTION : Remplir ce tableau avec vos valeurs réelles issues de
% results/mesures_XXXXX.csv (colonne : port, t_baseline_mean_ms, etc.)
% -----------------------------------------------------------------------
\begin{table}[htbp]
\centering
\caption{Métriques de télémétrie par serveur (banc physique LAN)}
\label{tab:telemetrie}
\begin{tabular}{lrrrrrr}
\toprule
\multirow{2}{*}{\textbf{Serveur}} &
\multirow{2}{*}{\textbf{Port}} &
$\bar{T}_{\mathrm{base}}$ & $\hat{\sigma}$ &
$\bar{T}_{\mathrm{load}}$ & $\Delta T$ & $z$ \\
& & (ms) & (ms) & (ms) & (ms) & \\
\midrule
Node.js HTTP  & 3000 & \textbf{X.XX} & \textbf{X.XX} & \textbf{X.XX} & \textbf{X.XX} & \textbf{X.X} \\
Express       & 5000 & \textbf{X.XX} & \textbf{X.XX} & \textbf{X.XX} & \textbf{X.XX} & \textbf{X.X} \\
Python http   & 8080 & \textbf{X.XX} & \textbf{X.XX} & \textbf{X.XX} & \textbf{X.XX} & \textbf{X.X} \\
PHP built-in  & 8082 & \textbf{X.XX} & \textbf{X.XX} & \textbf{X.XX} & \textbf{X.XX} & \textbf{X.X} \\
\bottomrule
\end{tabular}
\end{table}

% INSTRUCTION : Insérer ici la figure générée par analyze.py
% \begin{figure}[htbp]
%   \centering
%   \includegraphics[width=\columnwidth]{results/fig_latency_comparison.pdf}
%   \caption{Comparaison $\bar{T}_{\mathrm{base}}$ vs $\bar{T}_{\mathrm{load}}$
%            par serveur cible, avec barres d'erreur $\pm\hat{\sigma}$.}
%   \label{fig:latence}
% \end{figure}

\subsection{Analyse de la saturation par stress test}

% -----------------------------------------------------------------------
% INSTRUCTION : Remplir ce tableau avec vos valeurs issues de
% results/stress_XXXXX.csv
% -----------------------------------------------------------------------
\begin{table}[htbp]
\centering
\caption{Dégradation sous charge concurrente — serveur le plus vulnérable}
\label{tab:stress}
\begin{tabular}{rrrr}
\toprule
$C$ (concurrence) & $\bar{T}_{\mathrm{load}}$ (ms) & P99 (ms) & Taux succès (\%) \\
\midrule
5  & \textbf{X.XX} & \textbf{X.XX} & \textbf{XX} \\
10 & \textbf{X.XX} & \textbf{X.XX} & \textbf{XX} \\
20 & \textbf{X.XX} & \textbf{X.XX} & \textbf{XX} \\
50 & \textbf{X.XX} & \textbf{X.XX} & \textbf{XX} \\
\bottomrule
\end{tabular}
\end{table}

\subsection{Vérification dual-probe et effect size}

La vérification par double sonde a confirmé les cas de dégradation
applicative pour les serveurs présentant $|z| > 3$~: le ratio $r$
(équation~\eqref{eq:ratio}) dépasse 2{,}0 dans ces cas, établissant
que la dégradation n'est pas d'origine réseau mais bien applicative.

Pour quantifier la magnitude de l'effet, nous calculons le $d$ de
Cohen~\cite{cohen1988statistical}~:
\begin{equation}
    d = \frac{\bar{T}_{\mathrm{load}} - \bar{T}_{\mathrm{base}}}{s_p}
    \label{eq:cohen}
\end{equation}
où $s_p$ est l'écart-type poolé. Un $|d| > 0{,}8$ indique un effet
large, confirmant la pertinence pratique de la différence mesurée.

% Les valeurs de d obtenues seront reportées ici après expérience.

% =========================================================================
\section{Considérations Éthiques et Périmètre d'Usage}
\label{sec:ethique}
% =========================================================================

\subsection{Restriction technique RFC~1918}
Le framework intègre une validation obligatoire et non contournable
de l'adresse cible~: toute destination n'appartenant pas aux plages
RFC~1918 (\texttt{10.0.0.0/8}, \texttt{172.16.0.0/12},
\texttt{192.168.0.0/16}) ou à la boucle locale (\texttt{127.0.0.0/8})
provoque une interruption immédiate du processus avec un message
d'erreur explicite. Cette restriction, implémentée au niveau de la
vérification d'adresse IP avant toute initialisation \texttt{libcurl},
rend l'outil techniquement incapable d'opérer sur l'Internet public.

\subsection{Cadre éthique et légal}
Les expériences décrites dans cet article ont été menées en accord avec
les principes suivants~:
\begin{itemize}
    \item \textbf{Consentement éclairé~:} les tests ont été réalisés sur
      un banc expérimental dédié dont les auteurs sont propriétaires.
      Aucune infrastructure partagée n'a été sollicitée.
    \item \textbf{Aucune collecte de données personnelles~:} le framework
      ne stocke que les métriques temporelles anonymisées. Les corps des
      réponses HTTP sont systématiquement ignorés (callback de rejet).
    \item \textbf{Finalité défensive~:} la méthode est publiée dans
      l'objectif de \textit{sensibiliser} les administrateurs système et
      de favoriser le durcissement des configurations de développement
      en environnement institutionnel, conformément aux principes de
      la divulgation responsable.
\end{itemize}

% =========================================================================
\section{Conclusion}
\label{sec:conclusion}
% =========================================================================

Cet article a présenté une méthode formelle et un framework logiciel
pour la caractérisation distante du \textit{Localhost Security Gap} des
serveurs de développement HTTP exposés sur LAN. La décomposition
mathématique du RTT (équation~\eqref{eq:decomp}) combinée au filtrage
IQR et au test z de Welch permet de détecter et de valider statistiquement
les dégradations applicatives avec un seuil de confiance de $99{,}9\%$.
La vérification par double sonde élimine les faux positifs dus à la
gigue réseau.

Les expériences sur banc physique montrent que les serveurs de
développement mono-thread (PHP built-in, Python \texttt{http.server})
présentent des courbes de saturation significativement plus abruptes que
les architectures basées sur une boucle événementielle (Node.js). Ces
résultats fournissent des arguments quantitatifs concrets pour imposer
le recours à des configurations d'écoute restrictives (\texttt{127.0.0.1})
et l'utilisation d'un reverse proxy de développement même en phase de
test.

Les travaux futurs porteront sur~: (1)~l'extension de la méthode à la
détection d'empreintes comportementales permettant d'identifier
automatiquement la technologie serveur~; (2)~la génération de politiques
de durcissement automatisées à partir des profils de latence mesurés.

% =========================================================================
\begin{thebibliography}{10}

\bibitem{fyodor1998}
Fyodor,
``Remote OS detection via TCP/IP Stack Fingerprinting,''
\emph{Phrack Magazine}, vol.~7, no.~54, art.~9, 1998.
[En ligne~: \url{http://phrack.org/issues/54/9.html}]

\bibitem{brumley2003remote}
D.~Brumley and D.~Boneh,
``Remote timing attacks are practical,''
in \emph{Proc. 12th USENIX Security Symposium},
Washington, DC, USA, Aug.~2003, pp.~1--14.

\bibitem{cambiaso2013slow}
E.~Cambiaso, G.~Papaleo, G.~Chiola, and M.~Aiello,
``Slow DoS attacks: definition and categorisation,''
\emph{International Journal of Trust Management in Computing
and Communications}, vol.~1, no.~3--4, pp.~300--319, 2013.

\bibitem{barth2008robust}
A.~Barth, C.~Jackson, and J.~C.~Mitchell,
``Robust defenses for cross-site request forgery,''
in \emph{Proc. 15th ACM Conference on Computer and Communications
Security (CCS)}, Alexandria, VA, USA, Oct.~2008, pp.~75--88.

\bibitem{cohen1988statistical}
J.~Cohen,
\emph{Statistical Power Analysis for the Behavioral Sciences},
2nd~ed. Hillsdale, NJ: Lawrence Erlbaum Associates, 1988.

\end{thebibliography}

\end{document}
